\chapter{Source Code}

\begin{center}
\textbf{Originality Statement}\\[0.5em]
I verify that I am the sole author of the programs contained in this appendix, except where explicitly stated to the contrary.\\[1em]
\textit{Lian Matsuo}\\
\today
\end{center}

\vspace{1em}

The complete source code is submitted alongside this report. The main volume's Implementation chapter (Chapter~6) documents modelling and evaluation logic; this appendix is the canonical place for repository navigation, run paths, and artefact locations. The task-oriented map below highlights where to start for running sweeps, understanding model behaviour, and validating outputs.

%==============================================================================
\section{How to Navigate the Repository}
%==============================================================================

\textbf{Recommended reading path.}
\begin{enumerate}\tightlist
    \item Start with the top-level orchestration command in \code{scripts/run\_sweep.py}.
    \item Follow dataset defaults and sweep budgets in \code{src/lexisflow/config/datasets.py}.
    \item Read model backbones in \code{src/lexisflow/models/}.
    \item Inspect metrics in \code{src/lexisflow/evaluation/}.
    \item Validate assumptions with tests under \code{src/lexisflow/*/tests/}.
\end{enumerate}

\textbf{Core run path (end-to-end).}

\begin{center}
\small
\begin{tabularx}{\textwidth}{@{}>{\raggedright\arraybackslash}p{0.20\textwidth}>{\raggedright\arraybackslash}p{0.35\textwidth}X@{}}
\toprule
\textbf{Stage} & \textbf{Primary file(s)} & \textbf{Purpose} \\
\midrule
Sweep entrypoint & \code{scripts/run\_sweep.py} & Single command that orchestrates data prep, preprocessing, training, generation, and evaluation. \\
Dataset-specific drivers & \code{scripts/mimic/run\_sweep.py}, \code{scripts/challenge2012/run\_sweep.py} & Apply dataset profile defaults and execute the full $(n_t, n_{\text{noise}})$ grid. \\
Model training & \code{scripts/train\_hour0.py}, \code{scripts/train\_autoregressive.py} & Train IID hour-0 model and autoregressive transition model. \\
Evaluation and plots & \code{src/lexisflow/evaluation/*.py}, \code{scripts/common/analyze\_sweep.py} & Compute quality/privacy/TSTR metrics and generate sweep visualisations. \\
\bottomrule
\end{tabularx}
\normalsize
\end{center}

%==============================================================================
\section{High-Value Modules by Question}
%==============================================================================

\begin{center}
\small
\begin{tabularx}{\textwidth}{@{}>{\raggedright\arraybackslash}p{0.33\textwidth}>{\raggedright\arraybackslash}p{0.27\textwidth}X@{}}
\toprule
\textbf{If you need to understand...} & \textbf{Go to...} & \textbf{What to look for} \\
\midrule
How feature types are inferred & \code{src/lexisflow/data/feature\_utils.py} & Rules that classify columns as continuous, binary, or categorical before model fitting. \\
How autoregressive rows are built & \code{src/lexisflow/data/autoregressive.py} & Lag construction and static/dynamic feature separation for trajectory modelling. \\
How ForestFlow is implemented & \code{src/lexisflow/models/forest\_flow.py} & Conditional flow matching over discrete time levels with XGBoost regressors. \\
How HS3F differs from ForestFlow & \code{src/lexisflow/models/hs3f.py} & Heterogeneous sequential generation with separate handling of continuous and categorical targets. \\
How trajectory sampling works & \code{src/lexisflow/models/sampling.py} & Iterative autoregressive rollout of synthetic patient timelines. \\
How quality metrics are computed & \code{src/lexisflow/evaluation/quality\_metrics.py} & KS statistics, correlation preservation, and clinical range-violation calculations. \\
How privacy risk is measured & \code{src/lexisflow/evaluation/privacy\_metrics.py} & DCR-based summaries and DOMIAS-style membership inference implementation. \\
How downstream utility is measured & \code{src/lexisflow/evaluation/tstr\_framework.py} & Multi-task patient-level TSTR protocol for mortality and LOS prediction. \\
\bottomrule
\end{tabularx}
\normalsize
\end{center}

%==============================================================================
\section{Reproducibility Anchors}
%==============================================================================

\begin{center}
\small
\begin{tabularx}{\textwidth}{@{}>{\raggedright\arraybackslash}p{0.28\textwidth}>{\raggedright\arraybackslash}p{0.30\textwidth}X@{}}
\toprule
\textbf{Artifact type} & \textbf{Location} & \textbf{Why it matters} \\
\midrule
\textbf{Sweep defaults} & \code{src/lexisflow/config/datasets.py} & Canonical grid values, training-row budgets, and profile definitions used by both datasets. \\
\textbf{Primary results CSVs} & \code{results/sweep\_results.csv}, \code{results/challenge2012\_sweep\_results.csv} & Per-cell metrics and error fields used to generate plots and report tables. \\
\textbf{Preprocessor artifacts} & \code{artifacts/*preprocessor*.pkl} & Capture fitted transforms required for reproducible inverse-transform and generation. \\
\textbf{Model artifacts} & \code{artifacts/*/sweep/} & Persist trained models per grid cell for traceability and reruns. \\
\textbf{Unit tests} & \code{src/lexisflow/*/tests/} & Regression checks for data transforms, models, and evaluation metrics. \\
\textbf{Dependency lock/config} & \code{pyproject.toml}, \code{uv.lock} & Pin package versions and Python constraints for environment reproducibility. \\
\bottomrule
\end{tabularx}
\normalsize
\end{center}

\textbf{Note:} Complete source code listings are submitted electronically as a
companion archive. The program files used in the project examination are
identical to the files uploaded to KEATS.
